% Nouveau EGA V (cote 25): archive cote 25, site page 147 (PDF page 148).
% typescript. Draft compared: none.
% AI-transcribed by Claude Opus 5.5 (claude-opus-5-5, adaptive thinking, effort high), Message Batches API, batch ega25b (msgbatch_01A7n5YuXubUPNK1zsepi8RG), 2026-11-13; not yet checked by eye.
\scanpage{25/147}
\mspage{\hand{8}}

\uline{de} $\overset{\hand{\vee}}{\underline{P}}_{X/Y}\otimes\underline{O}_{V^1}$ , \uline{et en tant qu'homomorphisme dans c ette composante,}
\uline{il est déterminé précisément par lasedtion} $D = \text{dét}\,M(\hand{\phi})$ \uline{déjà envisagée}
\uline{par ailleurs} (tjrs modulo un signe à déterminer). Bien entendu,
l'homomorphisme b) s'en déduit, compte tenu de la commutativité énoncée
dans le lemme. Les points non supersinguliers$_{\text{x}}$ sont donc exacte\unsure{m}ent
ceux en lesquels $\underline{a}$: $\underline{\hand{\omega}}^{-2}_{X/Y}\otimes\underline{O}_{V^1}\hand{\longrightarrow}\underline{O}_{\unsure{V|X}^1}$ est un isomorphisme, et en
ces points,  le facteur $\underline{O}_{\unsure{V|X}^1}$ s'identif
fie donc exactement à l'intersection de $\underline{P}^1_{X/Y}\otimes\underline{O}_{V^1}$ et $\underline{\hand{\Omega}}^1_{Y/S}\otimes\underline{O}_{V^1}$
dans $\underline{\hand{\Omega}}^1_{X/S}\otimes\underline{O}_{V^1}$ . Donc sur l'ouvert de $V^1$ formé des points non
supersinguliers, on a une suite exacte
\[ 0\hand{\longrightarrow}\underline{O}_{V^1}\hand{\longrightarrow}\underline{\hand{\Omega}}^1_{Y/S}\otimes\underline{O}_{V^1}\hand{\longrightarrow}\underline{\hand{\Omega}}^1_{V^1/S}\hand{\longrightarrow} 0 \qquad , \]
ainsi que la suite exacte
\[ 0\hand{\longrightarrow}\underline{O}_{V^1}\hand{\longrightarrow}\underline{\hand{\Omega}}^1_{X/S}\otimes\underline{O}_{V^1}\hand{\longrightarrow}(\underline{\hand{\Omega}}^1_{X/Y}+\underline{\hand{\Omega}}^1_{Y/S})\otimes\underline{O}_{V^1}\hand{\longrightarrow} 0 \]

\placedmargin{$\|$} Il faudrait résumer  toute cette situation dans une proposition
récapitulative, et sous une forme légèrement plus générale, en partant
d'une section $\hand{\phi}$ d'un \uline{Module inversible}  K  X - ce qui oblige
à tordre par K un certain nombre des Modules introduits dans les
considérations précéde\unsure{nt|mm}es. C'est en effet  cette situation qu'on
rencontre  dans le par. des sections hyperplanes. Signaler aussi que
\placedmargin{|} $V(\hand{\phi})^{\text{sing}}$  et $V(\hand{\phi})^{\text{sup sing}}$ ne changent pas en r emplaçant $\hand{\phi}$ par u $\hand{\phi}$,
où u est une unité, \hand{en fait $d\hand{\phi}|V(\hand{\phi})$ et $M(\hand{\phi})$ sont multipliés par \unsure{$u$},}
et $D(\hand{\phi})$ est multiplié par $u^n$ ($n$ = dim. relative de $X$ sur $Y$),
\uline{Remarque} Sous les conditions envisagées sur X,Y,S,$\hand{\phi}$ , comme )
$U = V(\hand{\phi})^{\text{sing}}- V(\hand{\phi})^{\text{sup sing}}$ est non ramffié sur Y, ses fibres sur Y sont
discrètes, ce qui implique que pour tout $y\hand{\in} Y$, tout zéro singulier ordi-%
naire de $\hand{\phi}_y$ dans $X_y$ est isolé dans l'ensemble des zéros singuliers, 
i.e. est isolé dans l'ensemble des points non lisses de $V(\hand{\phi}_y)/k(y)$.
 On en conclut p.ex. aussitôt que si $d\hand{\geqslant} 2$, c'est un point géométri-%

\mspage{\hand{147}}
